% =====================================================================
%  Seksjon: Trenger vi den ytre droopen? To-droop-arkitekturen
%           (UIC-ens innebygde droop vs. den ytre WT-droopen).
%  Kilder:  src/tops_openfast/dyn_models/UIC.py       (perfect_tracking)
%           src/tops_openfast/dyn_models/windturbine_tower.py (_droop_command)
%  Bakgrunn: docs/STATUS_OG_VEIVALG.md, kap. 5
%  Preamble: amsmath, amssymb, graphicx, booktabs, siunitx
% =====================================================================

\section{Do We Need the Outer Droop? The Two-Droop Architecture}
\label{sec:outer-droop-necessity}

The frequency-support droop of \S\ref{sec:outer-droop} is not the only droop in
the model. The converter interface itself --- the Unified Integral Control (UIC)
grid-forming model --- \emph{is} a droop regulator by design, and it is active in
every run. It is therefore fair to ask whether the outer wind-turbine droop is
needed at all, or whether it merely duplicates what the converter already does.
This section separates the two mechanisms, explains what each contributes, and
states when the outer droop earns its place.

\subsection{The inner droop: UIC grid-forming control}
\label{subsec:odn-inner}

The UIC (\texttt{UIC\_sig}) implements the control law
\begin{equation}
  \dot{v}_i = j\,\omega_n K_i\,\varepsilon_s
            + \omega_n K_v\,\varepsilon_v
            + j\,v_i\,x_{\mathrm{filter}}\,\texttt{perfect\_tracking},
  \label{eq:odn-uic}
\end{equation}
where \(\varepsilon_s\) carries the active/reactive-power error and
\(\varepsilon_v\) the voltage error. The \(j\)-rotation turns a power error into a
frequency/angle adjustment (a P--f droop) and a voltage error into a magnitude
adjustment (a Q--V droop). This is an \emph{inherent, local, grid-forming} droop:
it gives the converter a synchronising, inertia-like response, and it is the very
reason a grid disturbance reaches the turbine at all rather than seeing a stiff
current source. The \(\texttt{perfect\_tracking}\) term \eqref{eq:odn-uic} can
cancel the droop and make the converter a PLL-free perfect tracker; with
\(\texttt{perfect\_tracking}=0\), as used throughout, the droop is \emph{on}.

Crucially, this inner droop is what makes the whole electro-mechanical study
possible: it is because the converter is droop-coupled, not stiff, that the
network disturbance propagates through \(P_e\) into \(T_e\) and hence into the
tower (\S\ref{sec:shaft-excitation}). Disabling it would decouple the turbine
from the grid and remove the interaction under study, which is why it is kept on
in every run.

\subsection{The outer droop: an explicit, de-loaded support service}
\label{subsec:odn-outer}

The outer droop \eqref{eq:od-delta} is a different object. It is not a
grid-forming primitive but an \emph{explicit dispatchable service} layered on the
power reference: it de-loads the turbine to hold a headroom reserve
\eqref{eq:od-base} and then releases that reserve in proportion to a measured
frequency deviation. The inner UIC droop does neither --- it responds to the
instantaneous power/voltage error at the terminal, but it does not curtail the
turbine to create a reserve, and it does not target a slow bus-frequency
excursion with a scheduled amount of extra energy. In power-system terms the
inner droop supplies the fast synchronising/inertial response, while the outer
droop supplies the slower primary frequency reserve.

\subsection{So: is it needed?}
\label{subsec:odn-verdict}

The honest answer is \emph{it depends on the question being asked}, and it is
worth stating plainly because it shaped the experimental design.

\begin{itemize}
  \item \textbf{For the electro-mechanical coupling to exist: no.} The grid
        reaches the tower through the inner UIC droop regardless of the outer
        loop. All the shaft-excitation and resonance results
        (\S\ref{sec:shaft-excitation}, \S\ref{subsec:nt-resonance}) are run with
        \(\texttt{droop\_enable}=0\); removing the outer droop does not change
        them. The torsional forced-response results of \S\ref{sec:forcedresponse}
        were likewise obtained with the outer droop off.
  \item \textbf{For demonstrating a realistic frequency-support trade-off: yes.}
        The value the outer droop adds is precisely that it lets the turbine
        \emph{act} on a frequency event with a controllable, de-loaded reserve,
        which is what exposes the central operational tension of this thesis: the
        same active-power injection that helps the grid frequency also drives the
        tower side-to-side mode harder. That trade-off (\S\ref{subsec:nt-tradeoff})
        cannot be posed without an explicit, toggleable support law, and the
        inner UIC droop --- always on, not de-loaded, not dispatchable --- does
        not provide it.
\end{itemize}

Two caveats qualify the outer droop and are carried into the results. First, it
reacts to a \emph{global} centre-of-inertia frequency \eqref{eq:od-fgrid}
computed in the driver from the generator speeds, rather than to a locally
measured bus frequency; this is a mild idealisation appropriate to a small,
coherent islanded grid but would need a local frequency estimator on a larger
network. Second, because the inner UIC droop is always active, the outer loop
sits \emph{on top of} an already frequency-responsive converter; the two are
complementary layers (fast local synchronising response plus slow dispatchable
reserve), not two copies of the same function. The recommendation, consistent
with the code as run, is therefore to keep the inner UIC droop on unconditionally
and to enable the outer droop only when a de-loaded primary-reserve service is
part of the scenario being studied.
